\documentclass[11pt]{article}

\usepackage{amsmath,amssymb,amsfonts,amsthm}
\usepackage{mathrsfs}
\usepackage{hyperref}
\usepackage{enumitem}
\usepackage{geometry}
\geometry{margin=1in}

\title{Gauge-Invariant Multiplicity Education System (MES 3.0)\\
A Defensive Publication and Technical Report}
\author{Citizen Gardens \\ The Foundation of Multiplicity}
\date{April 2026}

\begin{document}
\maketitle

\begin{abstract}
This document serves as a defensive publication and technical report formalizing the Gauge-Invariant Multiplicity Education System (MES 3.0). MES 3.0 integrates prime-based multiplicity algebra, gauge-equivalent curricular encodings, and a triad of degeneracy-lifting probes into a unified, empirically testable framework for self-correcting education systems. The core contributions are: (1) a gauge-invariant formulation of learner state spaces and curriculum factorizations, (2) a multi-scale dynamical system with explicit update rules for learners, curricular graphs, and prime sets, (3) a sparse, approximately commuting gauge map between alternative prime decompositions, and (4) a concrete family of degeneracy-lifting probes (Transfer, Perturbation, Representation Shift) instantiated for linear functions. The framework is designed to be falsifiable, operational, and extensible across domains.
\end{abstract}

\tableofcontents

\section{Executive Summary}

This report formalizes a line of development that unifies:
\begin{itemize}[itemsep=3pt]
  \item prime-based multiplicity algebra for encoding competencies;
  \item self-correcting, multi-scale educational dynamics;
  \item gauge symmetry across alternative curricular decompositions;
  \item degeneracy-lifting probes inspired by statistical and quantum mechanics.
\end{itemize}

\subsection{Problem Addressed}

Traditional adaptive learning systems treat:
\begin{itemize}[itemsep=3pt]
  \item the curriculum as a fixed graph of skills,
  \item learner knowledge as a vector of mastery scores,
  \item and errors as noise or individual deficits.
\end{itemize}

This constrains adaptation to local item routing and parameter tweaking, while leaving:
\begin{itemize}[itemsep=3pt]
  \item the underlying skill decomposition,
  \item its algebraic structure, and
  \item its consistency across institutions
\end{itemize}
largely unexamined and non-falsifiable.

\subsection{Core Idea}

MES 3.0 recasts an education system as a \emph{gauge-invariant multiplicity dynamical system}:
\begin{enumerate}[itemsep=3pt]
  \item A set of \emph{prime-labeled competencies} forms a basis for representing learner states as multiplicity vectors.
  \item The curriculum is a \emph{factorization space} rather than a static graph: composite skills correspond to multiplicative combinations of prime competencies.
  \item Learning dynamics are given by \emph{operators} acting on multiplicity vectors, with stochastic components.
  \item Crucially, there may be multiple valid prime bases (gauges) for the same domain; MES enforces \emph{gauge invariance} by learning sparse, approximately commuting maps between these bases.
  \item Degeneracy---distinct internal states with identical observable performance---is addressed via a \emph{probe library} (Transfer, Perturbation, Representation Shift) that lifts or confirms degeneracy in a controlled way.
\end{enumerate}

\subsection{Novelty}

Key novel elements:
\begin{itemize}[itemsep=3pt]
  \item \textbf{Gauge-equivalent prime sets:} multiple prime decompositions of the same curricular domain, connected by a learned gauge map, rather than a single arbitrary prime basis.
  \item \textbf{Gauge-invariant multiplicity:} learner multiplicity is defined as an equivalence class across gauges, not as coordinates in a privileged basis.
  \item \textbf{Dynamical gauge commutation:} the learned map $g_{12}$ is constrained by approximate commutation with learning operators, not just static representation alignment.
  \item \textbf{Degeneracy-lifting probes:} a triad of probes (near/far transfer, perturbation, representation shift) designed to distinguish true degeneracy from measurement failure.
  \item \textbf{Falsifiability:} the framework yields clear failure modes: non-existence of sparse commuting gauge maps or probe signatures that fail to predict out-of-distribution performance.
\end{itemize}

\subsection{Practical Deliverables}

This report provides:
\begin{itemize}[itemsep=3pt]
  \item A rigorous mathematical formulation of MES 3.0.
  \item A definition of gauge-equivalent prime systems and gauge maps.
  \item A concrete design of degeneracy-lifting probes for linear functions.
  \item Example code snippets for learning a sparse linear gauge map.
  \item A pilot study protocol for empirical validation.
\end{itemize}

\section{Mathematical Foundations of MES 3.0}

\subsection{Prime-Indexed Competency Spaces}

Let $\mathcal{D}$ be a curricular domain (e.g., introductory algebra). A \emph{prime decomposition} of $\mathcal{D}$ is a finite set of irreducible competencies
\[
P^{(i)} = \{p^{(i)}_1, \dots, p^{(i)}_{n_i}\},
\]
where each $p^{(i)}_k$ indexes a basic competency in decomposition $i$.

\begin{definition}[Learner State in a Fixed Gauge]
Given a prime set $P^{(i)}$, a learner state is a vector
\[
x^{(i)} \in \mathbb{N}^{n_i},
\]
where $x^{(i)}_k$ represents the multiplicity or robustness of competency $p^{(i)}_k$ for that learner (e.g., number of independent ways the competency can be reliably expressed).
\end{definition}

\begin{definition}[Composite Skill]
A composite skill $S$ is encoded as a vector $s^{(i)} \in \mathbb{N}^{n_i}$, interpreted as the prime factorization of $S$ in the $i$-th decomposition. The realization of $S$ by a learner with state $x^{(i)}$ is given by a threshold function
\[
f\left(x^{(i)}, s^{(i)}\right) \geq \theta,
\]
where $f$ could be, for example, a weighted inner product or a nonlinear capacity measure.
\end{definition}

\subsection{Multiplicity as an Equivalence Class}

Multiple decompositions can represent the same underlying learner state:

\begin{definition}[Gauge Family of Prime Sets]
A \emph{gauge family} for domain $\mathcal{D}$ is a finite collection
\[
\mathcal{G} = \{ P^{(1)}, P^{(2)}, \dots, P^{(k)} \},
\]
where each $P^{(i)}$ is a valid prime decomposition of $\mathcal{D}$.
\end{definition}

\begin{definition}[Observable Profile]
Let $\mathcal{F}$ be a family of tasks (e.g., assessments, probes). An observable profile for a learner is a vector
\[
I(L) \in \mathbb{R}^m,
\]
containing performance metrics (accuracy, latency, error type frequencies, probe responses) across $\mathcal{F}$.
\end{definition}

\begin{definition}[Gauge-Equivalent Representations]
Two learner representations $x^{(i)}$ and $x^{(j)}$ in gauges $i$ and $j$ are \emph{gauge-equivalent} if they produce indistinguishable observable profiles:
\[
x^{(i)} \sim x^{(j)} \quad \text{iff} \quad I^{(i)}(x^{(i)}) \approx I^{(j)}(x^{(j)}),
\]
for all observables in a specified family, up to experimental tolerance.
\end{definition}

\begin{definition}[Multiplicity as Equivalence Class]
The multiplicity of a learner is the equivalence class
\[
[x] = \{ x^{(i)} \mid i=1,\dots,k,\ x^{(i)} \sim x^{(j)} \ \forall j \},
\]
rather than any particular $x^{(i)}$. This is directly analogous to macrostates in statistical mechanics: many microstates (representations) per macro-observable profile.
\end{definition}

\subsection{Dynamical System Structure}

\subsubsection{Learner-Level Dynamics}

For a fixed gauge $i$, let $X_t^{(i)} \in \mathbb{N}^{n_i}$ denote a learner’s state at time $t$.

\begin{definition}[Learning Operators]
Each interaction type $\alpha$ (lesson, exercise, feedback event) is associated with an operator
\[
T^{(i)}_\alpha : \mathbb{N}^{n_i} \to \mathbb{N}^{n_i}.
\]
A stochastic evolution is given by
\[
X^{(i)}_{t+1} = T^{(i)}_{\alpha_t}(X^{(i)}_t) + \epsilon^{(i)}_t,
\]
where $\epsilon^{(i)}_t$ models noise (exploration, distraction, random fluctuations).
\end{definition}

\subsubsection{Error and Distortion Tensors}

Let $\hat{T}^{(i)}_\alpha$ be a learned prediction model (e.g., from historical data).

\begin{definition}[Error Tensor]
The \emph{error tensor} at time $t$ for interaction type $\alpha$ is
\[
\Delta^{(i)}_t = X^{(i)}_{t+1} - \hat{T}^{(i)}_\alpha(X^{(i)}_t) \in \mathbb{R}^{n_i}.
\]
This encodes component-wise distortions of multiplicity rather than a scalar error alone.
\end{definition}

\subsubsection{Multi-Scale Updates}

We define three time scales:
\[
\tau_{\text{learner}} \ll \tau_{\text{graph}} \ll \tau_{\text{prime}},
\]
for learner-level updates, curriculum graph updates, and prime refactorizations, respectively.

\paragraph{Local (Fast) Update.}
\begin{itemize}[itemsep=3pt]
  \item Learner state update: $x^{(i)} \leftarrow x^{(i)} + \eta \Delta^{(i)}$.
  \item Operator parameter update: adjust parameters of $\hat{T}^{(i)}_\alpha$ (e.g., via gradient descent) using $\Delta^{(i)}$.
\end{itemize}

\paragraph{Structural (Medium) Update.}

Let $A^{(i)}$ be the adjacency matrix of a competency graph in gauge $i$. When certain patterns in $\mathbb{E}[\Delta^{(i)}]$ recur above threshold, we modify $A^{(i)}$ (e.g., rewiring edges, adding or removing prerequisite relations).

\paragraph{Foundational (Slow) Update.}

When structural updates accumulate non-locally, we refactor the prime set:
\[
P^{(i)} \to P^{(i')}, \quad \phi^{(i\to i')}: \mathbb{N}^{n_i} \to \mathbb{N}^{n_{i'}},
\]
with a mapping $\phi$ chosen to preserve observable behavior as much as feasible.

\subsection{Gauge Maps and Dynamical Alignment}

\subsubsection{Gauge Map Definition}

\begin{definition}[Gauge Map]
A gauge map between two decompositions $P^{(1)}$ and $P^{(2)}$ is a function
\[
g_{12}: \mathbb{N}^{n_1} \to \mathbb{N}^{n_2},
\]
such that, for observed learner trajectories, the following hold approximately:
\begin{align*}
\text{(Representation alignment)} &:\quad g_{12}\left(X^{(1)}_t\right) \approx X^{(2)}_t, \\
\text{(Dynamical commutation)} &:\quad g_{12}\left(T^{(1)}_\alpha(x)\right) \approx T^{(2)}_\alpha\left(g_{12}(x)\right), \ \forall \alpha.
\end{align*}
\end{definition}

Typically $g_{12}$ is parameterized as a sparse linear map:
\[
g_{12}(x) = W x,\quad W \in \mathbb{R}^{n_2 \times n_1},\quad \|W\|_0 \text{ small}.
\]

\subsubsection{Learning the Gauge Map}

Consider stacked learner trajectories in each gauge:
\[
X_1 \in \mathbb{R}^{T \times n_1}, \quad X_2 \in \mathbb{R}^{T \times n_2}.
\]

\paragraph{Alignment Loss.}
\[
L_{\text{align}}(W) = \|X_2 - X_1 W^\top\|_F^2.
\]

\paragraph{Commutation Loss.}
Let $A^{(1)}_\alpha, A^{(2)}_\alpha$ be linear approximations to $T^{(1)}_\alpha, T^{(2)}_\alpha$. Then
\[
L_{\text{commute}}(W) = \sum_{t,\alpha} \left\| W A^{(1)}_\alpha X^{(1)}_t - A^{(2)}_\alpha W X^{(1)}_t \right\|^2.
\]

\paragraph{Regularized Objective.}
\[
\mathcal{L}(W) = L_{\text{align}}(W) + \mu L_{\text{commute}}(W) + \lambda \|W\|_1,
\]
where $\|W\|_1$ enforces sparsity and interpretability.

\subsection{Gauge-Invariant Multiplicity Theorem (Informal)}

\begin{quote}
\textbf{Gauge-Invariant Multiplicity Theorem (Informal).}
Multiplicity learning in a domain $\mathcal{D}$ is well-defined (i.e., multiplicity is an invariant of the learner, not the representation) if the following conditions hold:
\begin{enumerate}[itemsep=3pt]
  \item There exists a non-empty class of prime decompositions $\mathcal{G}$ for $\mathcal{D}$ such that for all $P^{(i)}, P^{(j)} \in \mathcal{G}$, a sparse gauge map $g_{ij}$ can be found.
  \item For each $g_{ij}$, the operator families $\{T^{(i)}_\alpha\}$ and $\{T^{(j)}_\alpha\}$ satisfy approximate commutation $g_{ij} \circ T^{(i)}_\alpha \approx T^{(j)}_\alpha \circ g_{ij}$ on empirical trajectories.
  \item Observable profiles $I(L)$ are invariant under basis change within $\mathcal{G}$ up to tolerance.
\end{enumerate}
When these conditions fail robustly (e.g., no sparse commuting map exists), multiplicity is representation-dependent in that domain.
\end{quote}

\section{Degeneracy and Probe Library}

\subsection{Degeneracy Notion}

\begin{definition}[Educational Degeneracy]
Two learners $L', L''$ are \emph{degenerate} at probe resolution $\beta$ if they exhibit indistinguishable observables for a specified task family:
\[
I_\beta(L') \approx I_\beta(L'').
\]
A system exhibits degeneracy when many distinct internal configurations lead to the same $I_\beta$.
\end{definition}

We want to distinguish:
\begin{itemize}[itemsep=3pt]
  \item \emph{True degeneracy}: internal differences do not matter for any plausible observable at a given scale.
  \item \emph{Measurement failure}: unmeasured internal structure has major consequences under richer observables (e.g., transfer, robustness).
\end{itemize}

\subsection{Probe Types}

For a fixed concept kernel $C$, we define three probe families:
\begin{enumerate}[itemsep=3pt]
  \item Transfer Probe $T_1$: near/far transfer (extensibility).
  \item Perturbation Probe $T_2$: noise and format shock (robustness).
  \item Representation Shift Probe $T_3$: non-trivial representational change (gauge flexibility).
\end{enumerate}

Each probe maps a macrostate equivalence class $[x]_{\text{base}}$ into higher-resolution observables, which may split it into sectors.

\section{Concrete Instantiation for Linear Functions}

We now instantiate the entire framework for the concept kernel of \emph{linear functions} in slope–intercept form.

\subsection{Concept Kernel}

\begin{itemize}[itemsep=3pt]
  \item Core relationship $y = mx + b$.
  \item Slope $m$ as rate of change and directional steepness.
  \item Intercept $b$ as starting value $y$ when $x = 0$.
  \item Interconversion between equation, graph, table, verbal description.
\end{itemize}

\subsection{Base Task Family $F_{\text{base}}$}

Define six items (5 multiple-choice, 1 open-response). Example types:
\begin{enumerate}[itemsep=3pt]
  \item Equation $\to$ value: given $y = 2x + 3$, find $y$ at $x=4$.
  \item Graph $\to$ slope: slope from two points on a line.
  \item Table $\to$ intercept: identify $y$-intercept from a table.
  \item Word problem $\to$ interpret slope: e.g., taxi fare, cost per mile.
  \item Graph $\to$ equation: matching slope and intercept visually.
  \item Open-response: write $y = -x + 4$ given slope and intercept.
\end{enumerate}

\paragraph{Scoring.}
\[
I_{\text{base}}(L) \in \{0,\dots,6\} \quad \text{(normalized to } [0,1]).
\]

\paragraph{Macrostate Definition.}
\[
[x]_{\text{base}} = \{ L \mid I_{\text{base}}(L) \geq \tau_{\text{base}} \},
\]
for a threshold $\tau_{\text{base}}$ (e.g., 5 or 6 out of 6), combined with near-stability constraints below.

\subsection{Near Transfer Family $F_{\text{near}}$}

Parallel items to $F_{\text{base}}$ with superficial changes (numbers, order, wording) but identical schema. Scoring $I_{\text{near}}$ is analogous.

\paragraph{Macrostate Stability Criterion.}
Learner $L$ is a stable member of $[x]_{\text{base}}$ if
\[
|I_{\text{near}}(L) - I_{\text{base}}(L)| \leq \delta_{\text{near}},
\]
for a small tolerance $\delta_{\text{near}}$ (e.g., 1 point).

\subsection{Far Transfer Family $F_{\text{far}}$ (Probe $T_1$)}

Tasks share the linear invariants but move to novel contexts:

\begin{itemize}[itemsep=3pt]
  \item Motion graphs: constant acceleration as slope of velocity–time graph.
  \item Recipe scaling: linear relation between quantity and servings.
  \item Unit conversion: Celsius–Fahrenheit conversion as linear function.
\end{itemize}

Score $I_{\text{far}}(L)$, and define:
\[
\Delta_1(L) = I_{\text{base}}(L) - I_{\text{far}}(L).
\]

\paragraph{Degeneracy Lifting.}
If the standard deviation of $\Delta_1$ in $[x]_{\text{base}}$ significantly exceeds test–retest noise, $T_1$ lifts degeneracy:
\[
\mathrm{FSI}_1 = \frac{\mathrm{SD}(\Delta_1)}{\mathrm{SD}_{\text{retest}}} > c,
\]
for some critical value $c$ (e.g., 1.96 for a 95\% confidence threshold).

\subsection{Perturbation Probe $T_2$}

\subsubsection{Channel $P_1$: Noise Injection}

Take $F_{\text{base}}$ items and add non-informational noise:
\begin{itemize}[itemsep=3pt]
  \item Extra numbers and irrelevant details in word problems.
  \item Visual clutter in graphs (extra lines, labels).
  \item Superfluous columns in tables.
\end{itemize}
Compute $I_{P_1}(L)$, define $\Delta_{1,\text{noise}}(L) = I_{\text{base}} - I_{P_1}$.

\subsubsection{Channel $P_2$: Format Shock}

Present structurally equivalent tasks in non-canonical formats:
\begin{itemize}[itemsep=3pt]
  \item Flowcharts representing $y = mx + b$.
  \item Sparse tables instead of graphs.
  \item Input–output machine descriptions of the same function.
\end{itemize}
Compute $I_{P_2}(L)$, define $\Delta_{1,\text{fmt}}(L)$ analogously.

Optionally aggregate:
\[
I_{\text{pert}}(L) = \frac{I_{P_1}(L) + I_{P_2}(L)}{2},\quad \Delta_2(L) = I_{\text{base}}(L) - I_{\text{pert}}(L).
\]

\subsection{Representation Shift Probe $T_3$}

Tightly linked to changes in internal mathematical representation:

\begin{itemize}[itemsep=3pt]
  \item Point-slope $\leftrightarrow$ slope–intercept conversion.
  \item Function notation: $f(x)$ instead of $y$.
  \item Recursive or parametric descriptions equivalent to $y=mx+b$.
\end{itemize}

Examples:
\begin{itemize}[itemsep=3pt]
  \item Given a line passing through $(1,5)$ with slope $2$, write its equation in $y=mx+b$ form.
  \item Let $f(x)=2x+3$, compute $f(4)$ and solve $f(a)=11$.
  \item “First mile costs 5 dollars, each additional mile costs 2 dollars” $\Rightarrow$ derive a formula equivalent to $y=2x+3$.
\end{itemize}

Define:
\[
I_{\text{rep}}(L),\quad \Delta_3(L) = I_{\text{base}}(L) - I_{\text{rep}}(L).
\]

\subsection{Sector Signatures}

For each learner in $[x]_{\text{base}}$, define a sector signature:
\[
\vec{\Delta}(L) = (\Delta_1(L), \Delta_2(L), \Delta_3(L)) \in \mathbb{R}^3.
\]

Clustering or thresholding on $\vec{\Delta}$ yields sectors. For example:

\begin{itemize}[itemsep=3pt]
  \item Conceptual-Flexible: all $\Delta_k$ small.
  \item Procedural-Brittle: all $\Delta_k$ large.
  \item Conceptual-Fragile: $\Delta_1 \approx 0$, $\Delta_2$ large, $\Delta_3 \approx 0$.
  \item Notation-Bound: moderate $\Delta_1$, small $\Delta_2$, large $\Delta_3$.
\end{itemize}

\section{Learning the Gauge Map: Example Code Snippet}

Below is an illustrative pseudocode-style snippet in Python-like syntax for learning a sparse linear gauge map $g_{12}(x) = W x$ between two prime gauges, using L1-regularized regression for alignment and a separate check for commutation.

\subsection{Data Preparation}

Assume:
\begin{itemize}[itemsep=3pt]
  \item $X_1$: array of shape $(T, n_1)$ with states in gauge 1.
  \item $X_2$: array of shape $(T, n_2)$ with corresponding states in gauge 2.
  \item $A1\_list$: list of operator matrices $A^{(1)}_\alpha$.
  \item $A2\_list$: list of operator matrices $A^{(2)}_\alpha$.
\end{itemize}

\subsection{Code Snippet}

\begin{verbatim}
import numpy as np
from sklearn.linear_model import Lasso

def fit_gauge_map(X1, X2, alpha=0.01):
    """
    Learn a sparse linear map W such that X2 ≈ X1 @ W^T.
    X1: (T, n1), X2: (T, n2)
    Returns W: (n2, n1)
    """
    T, n1 = X1.shape
    _, n2 = X2.shape
    W = np.zeros((n2, n1))

    # Fit each output dimension with L1-regularized regression
    for j in range(n2):
        y = X2[:, j]
        model = Lasso(alpha=alpha, fit_intercept=False)
        model.fit(X1, y)
        W[j, :] = model.coef_
    return W

def commutation_error(W, A1_list, A2_list, X1):
    """
    Compute mean squared commutation error:
    E[ || W A1 x - A2 W x ||^2 ] across samples and operators.
    """
    errors = []
    for A1, A2 in zip(A1_list, A2_list):
        # Apply operators in gauge 1 and map
        lhs = (X1 @ A1.T) @ W.T   # W A1 x
        # Map then apply operators in gauge 2
        rhs = (X1 @ W.T) @ A2.T   # A2 W x
        diff = lhs - rhs
        errors.append(np.mean(diff ** 2))
    return np.mean(errors)

# Example usage:
# W = fit_gauge_map(X1, X2, alpha=0.01)
# align_error = np.mean((X2 - X1 @ W.T)**2)
# comm_error = commutation_error(W, A1_list, A2_list, X1)
\end{verbatim}

The full objective combining alignment and commutation can be optimized iteratively:
\begin{itemize}[itemsep=3pt]
  \item fit $W$ for alignment;
  \item adjust operator matrices to reduce commutation error;
  \item iterate until convergence or diminishing returns.
\end{itemize}

\section{Pilot Study Protocol (Defensive Specification)}

\subsection{Hypotheses}

\begin{enumerate}[itemsep=3pt]
  \item There exist degenerate macrostates on linear functions under standard tests, i.e., many learners with high $I_{\text{base}}$ and $I_{\text{near}}$ but heterogeneous internal structure.
  \item The probe triad $(T_1, T_2, T_3)$ lifts this degeneracy into stable sectors characterized by $\vec{\Delta}$.
  \item Sector signatures predict out-of-distribution performance $I_{\text{OOD}}$ better than $I_{\text{base}}$ alone.
  \item In a dual-gauge setting, a sparse, approximately commuting gauge map exists that preserves sector structure across decompositions.
\end{enumerate}

\subsection{Participants}

\begin{itemize}[itemsep=3pt]
  \item $N \approx 120$--$150$ students enrolled in Algebra I.
  \item Inclusion: basic prior exposure to linear functions.
\end{itemize}

\subsection{Design}

\begin{itemize}[itemsep=3pt]
  \item Within-subject design: each learner completes $F_{\text{base}}$, $F_{\text{near}}$, $F_{\text{far}}$, T$_2$ items, and T$_3$ items.
  \item Two sessions to control fatigue, separated by 1--3 days.
  \item Randomized item order and interleaved filler items.
\end{itemize}

\subsection{Analysis Plan}

\begin{enumerate}[itemsep=3pt]
  \item Compute $I_{\text{base}}$, $I_{\text{near}}$, construct $[x]_{\text{base}}$.
  \item Compute $\Delta_1, \Delta_2, \Delta_3$ in $[x]_{\text{base}}$.
  \item Estimate test–retest SD using a subset who repeat base-like items.
  \item Check for degeneracy-lifting via FSI indices and multimodal structure in $\vec{\Delta}$.
  \item Cluster learners by $\vec{\Delta}$; interpret clusters as sectors.
  \item Administer $F_{\text{OOD}}$ and test whether $\vec{\Delta}$ improves prediction of $I_{\text{OOD}}$ beyond $I_{\text{base}}$.
  \item If dual-gauge data are collected, fit $W$ and evaluate alignment and commutation errors and sector preservation across gauges.
\end{enumerate}

\section{Discussion and Extensions}

\subsection{Theoretical Connections}

MES 3.0 sits at the intersection of:
\begin{itemize}[itemsep=3pt]
  \item Statistical mechanics: macrostates, microstates, multiplicity, and entropy.
  \item Quantum mechanics: degeneracy, spin multiplicity, selection rules, and degenerate perturbation theory.
  \item Representation learning: manifold alignment, linear probes, and transfer.
  \item Educational measurement: psychometrics, near/far transfer, robustness, and conceptual vs procedural knowledge.
\end{itemize}

\subsection{Failure Modes as Scientific Value}

The framework is deliberately fragile in the right places:
\begin{itemize}[itemsep=3pt]
  \item If no sparse, approximately commuting gauge map can be found, multiplicity is representation-dependent for that domain.
  \item If probe signatures fail to predict OOD performance, sector structure is likely an artifact.
  \item If sectors fail to stabilize across cohorts, the probe library is under-specified or the domain does not support clean sectors.
\end{itemize}

In all such cases, the failure is informative and refines the space of plausible multiplicity models.

\subsection{Future Work}

\begin{itemize}[itemsep=3pt]
  \item Extending the probe library to other kernels (quadratic functions, proportional reasoning, energy conservation).
  \item Incorporating affective and social dimensions into the multiplicity state (e.g., via tensor extensions).
  \item Exploring non-linear gauge maps and operator families (Koopman operator approaches).
  \item Integrating MES into live adaptive platforms to continuously learn $T_\alpha$, $g_{ij}$, and sector structure.
\end{itemize}

\section*{Conclusion}

This defensive publication captures a full architecture---mathematical, operational, and experimental---for a gauge-invariant multiplicity-based education system. It makes strong, testable claims about:
\begin{itemize}[itemsep=3pt]
  \item the nature of learner states as prime-indexed multiplicity vectors;
  \item the equivalence of different curricular prime decompositions under learned gauge maps;
  \item the role of carefully designed probes in lifting degeneracy and revealing hidden sectors;
  \item and the predictive power of these sectors for out-of-distribution performance.
\end{itemize}

The design is explicitly falsifiable: it will either reveal a robust, basis-independent structure underlying learning, or it will demonstrate the limits of multiplicity as a fundamental descriptor of educational dynamics.

\appendix
\section*{Mathematical Appendix}
\addcontentsline{toc}{section}{Mathematical Appendix}

\section{Preliminaries and Notation}

We recall the main objects of the Gauge-Invariant Multiplicity Education System (MES 3.0) and fix notation for the proofs.

\subsection{Spaces and Operators}

\begin{itemize}
  \item For each prime decomposition $P^{(i)} = \{p^{(i)}_1,\dots,p^{(i)}_{n_i}\}$, the learner state space is $X^{(i)} \equiv \mathbb{R}^{n_i}$ (relaxing to $\mathbb{R}$ for analysis).
  \item For each interaction type $\alpha$ (lesson, task, etc.), we have a linearized operator
  \[
    T^{(i)}_\alpha(x) = A^{(i)}_\alpha x + b^{(i)}_\alpha,
  \]
  with $A^{(i)}_\alpha \in \mathbb{R}^{n_i\times n_i}$ and $b^{(i)}_\alpha \in \mathbb{R}^{n_i}$.
  \item A gauge map between decompositions $i$ and $j$ is modeled as
  \[
    g_{ij}(x) = W_{ij} x, \quad W_{ij} \in \mathbb{R}^{n_j \times n_i},
  \]
  typically with sparsity constraints on $W_{ij}$.
\end{itemize}

We use $\|\cdot\|$ for the spectral (operator) norm on matrices and the Euclidean norm on vectors, and $\|\cdot\|_F$ for the Frobenius norm. When needed, we denote by $\|\cdot\|_1$ the elementwise $\ell_1$ norm on matrices or vectors.

\subsection{Approximate Commutation and Error Measures}

For a candidate gauge map $W_{ij}$ and operator families $\{A^{(i)}_\alpha\}$, $\{A^{(j)}_\alpha\}$, we define:

\begin{definition}[Pointwise Commutation Error]
For $x \in X^{(i)}$ and a fixed interaction type $\alpha$, the commutation residual is
\[
  r^{(ij)}_\alpha(x) := W_{ij} A^{(i)}_\alpha x - A^{(j)}_\alpha W_{ij} x.
\]
\end{definition}

\begin{definition}[Operator Commutation Error]
The (global) commutation error for $\alpha$ is
\[
  E^{(ij)}_\alpha(W_{ij}) := \sup_{\|x\|=1} \|r^{(ij)}_\alpha(x)\| = \| W_{ij} A^{(i)}_\alpha - A^{(j)}_\alpha W_{ij} \|.
\]
\end{definition}

For a finite family $\mathcal{A}$ of interaction types, we use the aggregate norm
\[
  E^{(ij)}(W_{ij}) := \max_{\alpha \in \mathcal{A}} E^{(ij)}_\alpha(W_{ij}).
\]

\section{Gauge Maps and Approximate Conjugacy}

This section makes precise the relationship between gauge maps and approximate dynamical conjugacy.

\subsection{Definition and Basic Properties}

\begin{definition}[Approximate Gauge Conjugacy]
Let $(X^{(i)}, \{A^{(i)}_\alpha\})$ and $(X^{(j)}, \{A^{(j)}_\alpha\})$ be two linear dynamical systems parametrized by interaction types $\alpha \in \mathcal{A}$. A linear map $W_{ij}: X^{(i)} \to X^{(j)}$ is an $\varepsilon$-\emph{gauge conjugacy} if
\[
  \| W_{ij} A^{(i)}_\alpha - A^{(j)}_\alpha W_{ij} \| \leq \varepsilon
\]
for all $\alpha \in \mathcal{A}$.
\end{definition}

Note that when $\varepsilon = 0$, this notion reduces to exact conjugacy between the linear operators.

\subsection{Stability of Trajectories under Gauge Maps}

We consider discrete-time trajectories under a fixed interaction type $\alpha$ in each gauge:
\begin{align*}
  x^{(i)}_{t+1} &= A^{(i)}_\alpha x^{(i)}_t, \\
  x^{(j)}_{t+1} &= A^{(j)}_\alpha x^{(j)}_t.
\end{align*}

\begin{theorem}[Trajectory Distortion Under Approximate Gauge Conjugacy]
\label{thm:traj-distortion}
Let $W_{ij}$ be an $\varepsilon$-gauge conjugacy between $(X^{(i)}, A^{(i)}_\alpha)$ and $(X^{(j)}, A^{(j)}_\alpha)$ for a fixed $\alpha$, and suppose
\[
  \|A^{(i)}_\alpha\| \leq \rho_i, \quad \|A^{(j)}_\alpha\| \leq \rho_j.
\]
Let $x^{(j)}_0 = W_{ij} x^{(i)}_0$. Then for all $t \ge 1$,
\[
  \left\| W_{ij} x^{(i)}_t - x^{(j)}_t \right\|
  \le t \, \varepsilon \, \|x^{(i)}_0\| \, \max(\rho_i,\rho_j)^{t-1}.
\]
\end{theorem}

\begin{proof}
We proceed by induction on $t$. For $t=0$, we have $W_{ij} x^{(i)}_0 = x^{(j)}_0$ by assumption, so the bound holds trivially.

Assume the bound holds for some $t \ge 0$; we consider $t+1$. We have
\[
  x^{(i)}_{t+1} = A^{(i)}_\alpha x^{(i)}_t, \quad
  x^{(j)}_{t+1} = A^{(j)}_\alpha x^{(j)}_t.
\]
Then
\begin{align*}
  W_{ij} x^{(i)}_{t+1} - x^{(j)}_{t+1}
  &= W_{ij} A^{(i)}_\alpha x^{(i)}_t - A^{(j)}_\alpha x^{(j)}_t \\
  &= \underbrace{\left( W_{ij} A^{(i)}_\alpha - A^{(j)}_\alpha W_{ij} \right) x^{(i)}_t}_{\text{commutation residual}} + A^{(j)}_\alpha \left( W_{ij} x^{(i)}_t - x^{(j)}_t \right).
\end{align*}
Taking norms:
\begin{align*}
  \left\| W_{ij} x^{(i)}_{t+1} - x^{(j)}_{t+1} \right\|
  &\le \|W_{ij} A^{(i)}_\alpha - A^{(j)}_\alpha W_{ij}\| \, \|x^{(i)}_t\|
     + \|A^{(j)}_\alpha\| \, \|W_{ij} x^{(i)}_t - x^{(j)}_t\| \\
  &\le \varepsilon \|x^{(i)}_t\| + \rho_j \, \|W_{ij} x^{(i)}_t - x^{(j)}_t\|.
\end{align*}
We also have by repeated application of the operator norm bound:
\[
  \|x^{(i)}_t\| \le \rho_i^t \|x^{(i)}_0\|.
\]
Thus
\begin{align*}
  \left\| W_{ij} x^{(i)}_{t+1} - x^{(j)}_{t+1} \right\|
  &\le \varepsilon \rho_i^t \|x^{(i)}_0\| + \rho_j \cdot t \varepsilon \|x^{(i)}_0\| \max(\rho_i,\rho_j)^{t-1} \\
  &\le \varepsilon \|x^{(i)}_0\| \max(\rho_i,\rho_j)^t + t \varepsilon \|x^{(i)}_0\| \max(\rho_i,\rho_j)^t \\
  &= (t+1) \varepsilon \|x^{(i)}_0\| \max(\rho_i,\rho_j)^t,
\end{align*}
which gives the desired bound for $t+1$. The base case completes the induction.
\end{proof}

\begin{corollary}[Uniform Boundedness Under Contraction]
\label{cor:bounded-contraction}
If additionally $\max(\rho_i,\rho_j) \le 1$, then
\[
  \| W_{ij} x^{(i)}_t - x^{(j)}_t \| \le t\,\varepsilon \,\|x^{(i)}_0\|.
\]
If $\max(\rho_i,\rho_j) < 1$, then the distortion grows at most linearly for small $t$ and is uniformly bounded for all $t$.
\end{corollary}

\section{Probe Operators and Sector Separation}

We formalize the idea that degeneracy-lifting probes act as ``selection-rule'' operators in the observable space.

\subsection{Probe Operators as Maps on Observable Profiles}

Let $\mathcal{F}_{\text{base}}$ be the base task family and let an observable profile be $I_{\text{base}}: X \to \mathbb{R}$ mapping a state to performance on the base tasks. For a probe family $\beta \in \{1,2,3\}$, define $I_\beta: X \to \mathbb{R}$ as performance on that probe family.

\begin{definition}[Probe-Induced Degeneracy Map]
For a given equivalence class
\[
  [x]_{\text{base}} := \{x' \in X \mid I_{\text{base}}(x') \approx I_{\text{base}}(x)\},
\]
we define the probe-induced map
\[
  \Phi_\beta([x]_{\text{base}}) := \{ I_\beta(x') \mid x' \in [x]_{\text{base}} \}.
\]
\end{definition}

A probe \emph{lifts degeneracy} at resolution $\beta$ if the image $\Phi_\beta([x]_{\text{base}})$ has non-trivial spread beyond measurement noise.

\subsection{Sector Separation Criterion}

Assume that for a fixed concept kernel $C$, the probe family produces scores $I_{\text{base}}, I_1, I_2, I_3$ and corresponding drops
\[
  \Delta_k(x) = I_{\text{base}}(x) - I_k(x), \quad k=1,2,3.
\]

\begin{definition}[Sector Signature Map]
Define the mapping
\[
  \Sigma: X \to \mathbb{R}^3, \quad \Sigma(x) := (\Delta_1(x), \Delta_2(x), \Delta_3(x)).
\]
For the macrostate $[x]_{\text{base}}$, let
\[
  S_{\text{base}} = \{ \Sigma(x') \mid x' \in [x]_{\text{base}} \}.
\]
\end{definition}

\begin{proposition}[Necessary Condition for Sector Non-Triviality]
\label{prop:sector-nontrivial}
If $S_{\text{base}}$ is contained in a ball of radius $r$ in $\mathbb{R}^3$, with $r$ less than or equal to the empirically estimated test--retest noise radius $r_{\text{noise}}$, then no non-trivial sector decomposition can be inferred from $(T_1,T_2,T_3)$ alone at the chosen resolution.
\end{proposition}

\begin{proof}
If $S_{\text{base}} \subset B(0,r)$ with $r \le r_{\text{noise}}$, then all observed variation in $\Sigma(x)$ is consistent with measurement noise. Any clustering of $S_{\text{base}}$ would be unstable under re-measurement and thus cannot be distinguished from random fluctuation. Hence no statistically justifiable sector boundaries can be drawn from these data alone.
\end{proof}

\begin{theorem}[Cluster Separation Under Probe Triad]
\label{thm:cluster-separation}
Assume there exist subsets $U_1,\dots,U_m \subset [x]_{\text{base}}$ such that:
\begin{enumerate}
  \item For each $\ell$, $\Sigma(U_\ell)$ is contained in a ball of radius $r_{\text{within}}$.
  \item For any $\ell \ne k$, the distance between the centers $\mu_\ell, \mu_k$ of $\Sigma(U_\ell)$ and $\Sigma(U_k)$ satisfies
  \[
    \|\mu_\ell - \mu_k\| \ge d_{\text{between}} > 2 r_{\text{within}} + 2 r_{\text{noise}}.
  \]
\end{enumerate}
Then there exists a clustering of $S_{\text{base}}$ into $m$ sectors such that, with high probability under repeated measurement, each sector remains stable and distinct.
\end{theorem}

\begin{proof}
By assumption, any element of $\Sigma(U_\ell)$ lies within $r_{\text{within}}$ of $\mu_\ell$, and measurement noise may shift any $\Sigma(x)$ by at most $r_{\text{noise}}$ (in radius, with high probability). Therefore, the effective radius of uncertainty around $\mu_\ell$ is $r_{\text{within}} + r_{\text{noise}}$.

Since $d_{\text{between}} > 2(r_{\text{within}}+r_{\text{noise}})$, the uncertainty regions around different cluster centers are disjoint. This implies that, under repeated measurement, points associated with $U_\ell$ will not cross into the uncertainty region of another cluster $U_k$. Thus we can assign sector labels by nearest-center classification, and these labels will be preserved under noise with high probability.
\end{proof}

\section{Multi-Scale Stability of MES 3.0}

We now formalize the separation of time scales and show a simple stability result.

\subsection{Time Scale Separation}

Let $\mathcal{T}_L$, $\mathcal{T}_G$, $\mathcal{T}_P$ be the sets of time indices for learner-level, graph-level, and prime-level updates, respectively, with:

\[
  \mathcal{T}_L \supset \mathcal{T}_G \supset \mathcal{T}_P,
\]
and characteristic step sizes
\[
  |\mathcal{T}_L| \gg |\mathcal{T}_G| \gg |\mathcal{T}_P|
\]
over any shared horizon.

At learner times $t \in \mathcal{T}_L \setminus \mathcal{T}_G$:
\[
  X_{t+1} = T_{\alpha_t}(X_t) + \epsilon_t,
\]
with fixed $A_\alpha$, $P$, and graph $A$.

At graph times $t \in \mathcal{T}_G \setminus \mathcal{T}_P$:
\[
  A_{t+1} = A_t + \Gamma(\mathbb{E}[\Delta_t]),
\]
with fixed $P$.

At prime times $t \in \mathcal{T}_P$:
\[
  P \to P', \quad X \to \phi(X),
\]
with a refactorization mapping $\phi$.

\subsection{Stability Under Small Operator Residuals}

\begin{definition}[MES Stability]
The system is \emph{stable} over horizon $[0,T]$ if:
\begin{enumerate}
  \item All operator residuals remain bounded:
  \[
    \sup_{t \le T} \sup_\alpha \|r_t(\alpha)\| \le R < \infty.
  \]
  \item The number of graph updates and the magnitude of each update yield a bounded total graph variation:
  \[
    \sum_{t \in \mathcal{T}_G \cap [0,T]} \|A_{t+1} - A_t\| \le V_G < \infty.
  \]
  \item The number of prime refactorizations in $[0,T]$ is finite, and the associated maps $\phi$ are uniformly Lipschitz.
\end{enumerate}
\end{definition}

\begin{theorem}[Sufficient Condition for Bounded Learner Trajectories]
Assume:
\begin{enumerate}
  \item All learner-level operators $T_{\alpha_t}$ are uniformly bounded: $\|T_{\alpha_t}\| \le \rho$ with $\rho < \infty$.
  \item The noise terms $\epsilon_t$ are bounded: $\|\epsilon_t\| \le \sigma$ for all $t$.
  \item MES is stable over $[0,T]$ in the above sense.
\end{enumerate}
Then for any learner initial condition $X_0$, the trajectory $\{X_t\}_{t=0}^T$ remains bounded:
\[
  \sup_{t \le T} \|X_t\| \le B,
\]
for some $B$ depending on $\|X_0\|$, $\rho$, $\sigma$, $R$, $V_G$, and the Lipschitz constants of the prime refactorization maps.
\end{theorem}

\begin{proof}[Sketch of Proof]
Between graph and prime updates, the system is governed by a fixed family of bounded operators with bounded noise, hence standard discrete-time stability bounds apply:
\[
  \|X_{t+1}\| \le \rho \|X_t\| + \sigma + R.
\]
By iterating this inequality and using $\rho < \infty$, we obtain a bound over any subinterval between structural changes.

Graph updates induce bounded perturbations of the operators due to finite total variation $V_G$, which can be absorbed into the constants appearing in the bound (via standard perturbation arguments in operator theory). Prime refactorizations, being finitely many and Lipschitz, transform the state space in a way that scales the bound by at most a finite product of Lipschitz constants. Combining these pieces yields a global finite bound $B$.
\end{proof}

\section{Remarks on Extensions}

\subsection{Nonlinear Operators and Koopman Lifts}

The above analysis uses linearized operators. In practice, learning operators may be nonlinear:
\[
  X_{t+1} = F_\alpha(X_t) + \epsilon_t.
\]
A standard extension is to consider a Koopman operator $K_\alpha$ acting on an embedding of observables $\psi(X)$, with an approximate linear structure:
\[
  \psi(X_{t+1}) \approx K_\alpha \psi(X_t).
\]
Gauge maps and probe analyses can then be carried out in the lifted space with analogous operator norm bounds.

\subsection{Entropy of Gauge Classes}

Given a macrostate equivalence class $[x]$ and a finite discretization of state space, one can define a multiplicity count $\Omega([x])$ as the cardinality of compatible microstates (representations) across gauges. An entropy-like quantity
\[
  S([x]) = \log \Omega([x])
\]
then quantifies the degeneracy of the macrostate. Probe-induced sector splits correspond to reductions in $\Omega$ and $S$; formalizing this requires a combinatorial model of the state space and gauge family, which is left for future work.

\section*{Conclusion}

This appendix has provided:
\begin{itemize}
  \item operator norm bounds for trajectory distortion under approximate gauge conjugacy;
  \item a formal sector separation criterion based on probe-induced signatures;
  \item a multi-scale stability result for MES 3.0 under bounded residuals and finite structural variation.
\end{itemize}
These results close the loop between the conceptual MES framework and explicit, analyzable mathematical properties, strengthening its status as a candidate empirical theory of learning dynamics rather than a purely metaphorical construction.


\nocite{*}
\bibliographystyle{plainnat}
\bibliography{references}
\end{document}